\answer{ex:11:1}{(а)~$T=k^2/\bigl(\Phi(\Delta I/I)^2\bigr)=9\,\text{с}$. (б)~$45$ датчиків, пляма $\approx7\times7$, роздільність гірша в $\approx7$ разів.}
\answer{ex:11:2}{(а)~$2$--$3.3$ біта на імпульс, $22$--$34\,\%$ межі ($9.1$--$9.8$ біта). (б)~$\log_2\binom{400}{20}\approx111$ бітів; у середньому $1$ спільний тип.}
\answer{ex:11:3}{(а)~$\sigma/9\le I\le9\sigma$, $1.9$ порядку. (б)~$6$ і $7$.}
\answer{ex:11:4}{(а)~$f<343/0.44\approx780$~Гц. (б)~$625$ імпульсів, $125$ нейронів.}
\answer{ex:11:5}{$21$ біт на подію, $4.8\cdot10^5$ подій/с, $7$ подій на піксель краю: $136$ пікселів/с. Звичайна камера --- $1.25$ кадру/с.}
\answer{ex:11:6}{У логарифмах $\approx1.0\,\text{с}$ в обидва боки (теорія $0.96\tau$); лінійно $0.2\,\text{с}$ угору і $3.0\,\text{с}$ униз (теорія $0.18\tau$ і $3.0\tau$).}
\answer{ex:11:7}{(а)~$\ln I$ розподілений логістично з медіаною $\ln\sigma$ і масштабом $1$: розкид $\pi/\sqrt3$ у натуральних логарифмах, $0.79$ порядку. (б)~$r=r_{\max}\Phi_I(I)^2$.}
